\documentclass[11pt]{article}

% Packages for XeLaTeX
\usepackage{fontspec}
\usepackage{amsmath,amssymb,amsthm}
\usepackage{geometry}
\usepackage{hyperref}
\usepackage{enumitem}
\usepackage{mathtools}
\usepackage{bm}
\usepackage{framed}
\usepackage{xcolor}



\geometry{margin=1in}

% Custom commands
\newcommand{\E}{\mathbb{E}}
\newcommand{\R}{\mathbb{R}}
\newcommand{\N}{\mathcal{N}}
\newcommand{\ELBO}{\text{ELBO}}
\newcommand{\tr}{\text{tr}}
\setlength{\parindent}{0pt}

\title{Project Planning Report:\\
Optimisation Methods for Variational Inference}
\author{David Ewing (82171165)\\
ENEL445  \\
University of Canterbury}
\date{17 March 2026}

\begin{document}

\maketitle

\section*{Background and Problem Intuition}

 

Bayesian inference provides principled uncertainty quantification but requires computing intractable posterior distributions. For most real-world models, the posterior integral, \par

\begin{equation*}
    p(\bm{\theta} \mid \bm{y}) = \frac{p(\bm{y} \mid \bm{\theta})p(\bm{\theta})}{p(\bm{y})},
\end{equation*}
\noindent 
cannot be evaluated analytically. Variational Inference (VI) converts this inference problem into an optimisation problem, making it directly amenable to the gradient-based methods in this course. \vspace{2em}

 

Instead of computing the true posterior, VI finds the ``best'' approximation $q(\bm{\theta})$ from a tractable family by maximising the Evidence Lower Bound (ELBO). This transforms intractable integration into tractable optimisation---I think this is a natural fit for an optimisation course project.

\section*{Key Notation}

\begin{center}
\begin{tabular}{ll}
\hline
Symbol & Meaning \\
\hline
$\bm{\beta}$ & Regression coefficient vector \\
$\bm{u}$ & Random effects vector (hierarchical model) \\
$\bm{\varepsilon}$ & Observation noise vector \\
$\tau_e$ & Observation-noise precision \\
$\tau_u$ & Random-effects precision \\
$\bm{\mu}_\beta$ & Mean of variational posterior $q(\bm{\beta})$ \\
$\bm{\Sigma}_\beta$ & Covariance of variational posterior $q(\bm{\beta})$ \\
\hline
\end{tabular}
\end{center}
\newpage
\section*{Optimisation Problem Formulation}

\subsection*{Decision Variables}

For Bayesian linear regression with $p$ predictors, the variational parameters (decision variables) are:
\begin{itemize}
\item $\bm{\mu}_\beta \in \R^p$: posterior mean of regression coefficients
\item $\bm{\Sigma}_\beta \in \R^{p \times p}$: posterior covariance (symmetric positive definite)
\item $a_e, b_e > 0$: shape and rate parameters of the variational Gamma factor for precision
\end{itemize}

Here, $a_e$ and $b_e$ parameterise $q(\tau_e)=\mathrm{Gamma}(a_e,b_e)$; the model prior on $\tau_e$ is specified in the Test Case section.

Total dimension: $d = p + \frac{p(p+1)}{2} + 2 = 11$ variables for $p=3$ predictors.

\subsection*{Objective Function}

Maximise the Evidence Lower Bound (ELBO):

\begin{align}
\ELBO(\bm{\phi})
&= \E_q[\log p(\bm{y}, \bm{\beta}, \tau_e)] + H(q) \\
&= \E_q[\log p(\bm{y}, \bm{\beta}, \tau_e)] - \E_q[\log q(\bm{\beta}, \tau_e)] .
\end{align}


where $\bm{\phi} = (\bm{\mu}_\beta, \bm{\Sigma}_\beta, a_e, b_e)$ are the variational parameters.

\subsection*{Constraints}

\begin{itemize}
\item $\bm{\Sigma}_\beta \succ 0$: Covariance matrix must be positive definite
\item $a_e, b_e > 0$: Gamma parameters must be positive
\item Optional: Minimum variance constraints $\sigma^2_{\beta_j} \geq \sigma^2_{\text{min}}$ to prevent underdispersion
\end{itemize}

These constraints can be handled via parameterisation (e.g., Cholesky decomposition for $\bm{\Sigma}_\beta$) or penalty methods.

\subsection*{Computational Characteristics}

\begin{itemize}
\item Gradient computation: $O(np^2)$ per evaluation (dominated by matrix operations)
\item Hessian computation: $O(np^2 + d^3)$ for Newton's method
\item Parameters exhibit coupling (changing $\bm{\mu}_\beta$ affects optimal $\bm{\Sigma}_\beta$)
\item Must handle special functions: digamma $\psi(\cdot)$, trigamma $\psi'(\cdot)$ for Gamma distribution
\end{itemize}
\newpage
\section*{Test Case: Bayesian Linear Regression}

The test case uses Bayesian linear regression with:
\begin{align}
\bm{y} &\sim \N(\bm{X}\bm{\beta}, \tau_e^{-1} \bm{I}_n) \quad \text{(likelihood)}\\
\bm{\beta} &\sim \N(\bm{0}, \lambda_\beta^{-1} \bm{I}_p) \quad \text{(prior on coefficients)}\\
\tau_e &\sim \text{Gamma}(\alpha_e, \gamma_e) \quad \text{(prior on precision)}
\end{align}

Mean-field approximation:
\begin{align}
q(\bm{\beta}, \tau_e) &= q(\bm{\beta}) \, q(\tau_e)\\
q(\bm{\beta}) &= \N(\bm{\mu}_\beta, \bm{\Sigma}_\beta)\\
q(\tau_e) &= \text{Gamma}(a_e, b_e)
\end{align}

Variational parameters: $\bm{\phi} = (\bm{\mu}_\beta, \bm{\Sigma}_\beta, a_e, b_e)$ with dimension $d = p + \frac{p(p+1)}{2} + 2$

For $p=3$ predictors: $d = 3 + 6 + 2 = 11$ parameters to optimise.
 
\section*{Overview of Optimisation Methods}

\subsection*{Methods considered}

\begin{itemize}
\item CAVI: Domain-specific baseline exploiting problem structure
\item Gradient Ascent: First-order baseline method using gradient information and line search
\item Newton's Method: Second-order method with quadratic local convergence near a well-conditioned optimum
\item BFGS: Quasi-Newton method that approximates Hessian information; practical general-purpose method for smooth unconstrained problems
\end{itemize}

\subsection*{Baseline Method: CAVI\footnote{Blei, Kucukelbir, and McAuliffe (2017), Variational Inference: A Review for Statisticians (Journal of the American Statistical Association).}}

Coordinate Ascent Variational Inference (CAVI) is VI-specific. It updates each variational parameter block sequentially using closed-form solutions:

\begin{align}
\bm{\Sigma}_\beta &\leftarrow \left(\frac{a_e}{b_e} \bm{X}^T\bm{X}\right)^{-1}\\
\bm{\mu}_\beta &\leftarrow \frac{a_e}{b_e} \bm{\Sigma}_\beta \bm{X}^T\bm{y}\\
a_e &\leftarrow \alpha_e + \frac{n}{2}\\
b_e &\leftarrow \gamma_e + \frac{1}{2}\left[\|\bm{y} - \bm{X}\bm{\mu}_\beta\|^2 + \tr(\bm{X}^T\bm{X}\bm{\Sigma}_\beta)\right]
\end{align}

Properties: Monotonic ELBO increase \footnote{Monotonic ELBO increase means the objective is non-decreasing across iterations: $\mathcal{L}^{(t+1)} \ge \mathcal{L}^{(t)}$. In CAVI, each coordinate update maximises the ELBO with respect to one variational factor while holding the others fixed, so each step cannot decrease the ELBO.}, linear convergence, $O(np^2)$ per iteration. Provides analytical benchmark for validation.


\newpage 

\subsection*{Mapping ENEL445 Methods to VI Problem}

\begin{enumerate}[leftmargin=*]
\item Parameter Ascent (Course: gradient descent)
\begin{itemize}
\item Standard iterative update: $\bm{\phi}^{(t+1)} = \bm{\phi}^{(t)} + \alpha_t \nabla \ELBO(\bm{\phi}^{(t)})$
\item Line search from Chapter 3.2 (bisection/Wolfe conditions)
\item Stopping criteria: $\|\nabla \ELBO\| < \epsilon$
\item Convergence characteristics: assuming it follows a linear rate
\end{itemize}

\item Newton's Method   
\begin{itemize}
\item Standard Newton update: $\bm{\phi}^{(t+1)} = \bm{\phi}^{(t)} + \alpha_t [\nabla^2 \ELBO]^{-1} \nabla \ELBO$
\item Requires computing and inverting the Hessian matrix
\item Computational cost: $O(d^3)$ per iteration for Hessian inversion
\item Convergence characteristics: unknown. 
\item Challenge: Hessian may not be positive definite away from the optimum
\end{itemize}

\item BFGS Quasi-Newton (one of five required methods from Project 2)
\begin{itemize}
\item BFGS update formula builds inverse Hessian approximation
\item Uses gradient differences: $\bm{s}_k = \bm{\phi}^{(k)} - \bm{\phi}^{(k-1)}$, $\bm{y}_k = \nabla \ELBO(\bm{\phi}^{(k)}) - \nabla \ELBO(\bm{\phi}^{(k-1)})$ 
\item Computational cost: $O(d^2)$ per iteration
\item Convergence characteristics: unknown.  
\end{itemize}
\end{enumerate}


Comparison will reveal whether general optimisation methods can compete with or outperform problem-specific algorithms

\section*{Solution Plan}

\subsection*{Phase 1: Data Generation and Baseline}

Data simulation I (complete):
\begin{itemize}
\item Generate $n = 100$ observations, $p = 3$ predictors
\item Design matrix $\bm{X} \sim \N(0, 1)$, true coefficients $\bm{\beta}_{\text{true}} = [1, -0.5, 2]^T$
\item Noise: $\epsilon \sim \N(0, 0.5^2)$
\end{itemize}

Mathematical foundations (in progress):
\begin{itemize}
\item Complete ELBO derivation for Bayesian linear regression
\item CAVI closed-form update equations
\item Gradient expressions: $\nabla_{\bm{\mu}_\beta} \ELBO$, $\nabla_{\bm{\Sigma}_\beta} \ELBO$, $\nabla_{a_e} \ELBO$, $\nabla_{b_e} \ELBO$
\item Hessian block structure: diagonal and mixed blocks
\item BFGS update formulation with Woodbury identity
\end{itemize}

Derivations to be documented in \texttt{vi\_optimisation\_solution.pdf} (or an appendix to the report?)

\subsection*{Phase 2: Implementation (in progress)}

\begin{enumerate}[leftmargin=*]
\item Data Generation (5\%)
\begin{itemize}
\item Simulate data: $n = 100$ observations, $p = 3$ predictors
\item Generate $\bm{X}$ from standard normal, $\bm{\beta}_{\text{true}} = [1, -0.5, 2]^T$
\item Add noise: $\epsilon \sim \N(0, 0.5^2)$
\end{itemize}

\item Baseline Implementation: CAVI (15\%)
\begin{itemize}
\item Implement closed-form coordinate updates
\item Track ELBO convergence
\item Stopping criterion: $|\ELBO^{(t+1)} - \ELBO^{(t)}| < 10^{-6}$
\end{itemize}

\item Method 1: Gradient Ascent (15\%)
\begin{itemize}
\item Implement ELBO gradient computation
\item Line search: bisection method or Wolfe conditions
\item Stopping criterion: $\|\nabla \ELBO\| < 10^{-6}$
\item Track iteration count and ELBO values
\end{itemize}

\item Method 2: Newton's Method (0\%)
\begin{itemize}
\item Implement Hessian computation (block structure)
\item Regularisation: can I add $\lambda \bm{I}$ if Hessian not positive definite? 
\item Line search for step size $\alpha_t$
\item Track convergence and computational cost
\end{itemize}

\item Method 3: BFGS Quasi-Newton (0\%)
\begin{itemize}
\item Initialise inverse Hessian approximation: $\bm{B}_0 = \bm{I}_d$ 
\item Track superlinear convergence
\end{itemize}

\item Comparison and Analysis (5\%)
\begin{itemize}
\item Compare convergence speed (iterations to $10^{-6}$ tolerance)
\item Compare computational cost (time per iteration, total time)
\item Compare solution quality (posterior approximation accuracy)
\item Analyse when each method performs best
\end{itemize}
\end{enumerate}
\newpage
\subsection*{Success Metrics}

\begin{enumerate}
\item Convergence: All methods reach ELBO optimum within tolerance
\item Correctness:
\begin{itemize}
\item CAVI correctness: Gradient/Newton/BFGS implementation is correct if it reaches the same ELBO optimum / variational fixed point as CAVI (same VI objective, same variational family)
\item MCMC correctness: VI is assessed against MCMC as a posterior-accuracy benchmark (means, intervals, calibration), not expected to match exactly because VI is approximate
\end{itemize}
\item Performance hierarchy:
\begin{itemize} 
\item CAVI likely fastest for this conjugate problem
\item MCMC/Gibbs expected slowest computationally, but best for posterior uncertainty accuracy (gold-standard benchmark, not a direct ELBO-optimiser competitor)
\end{itemize}
 
\end{enumerate}
 
\section*{Progress To Date}

\subsection*{Work in progress using an iterative / Agile method}

\begin{enumerate}
\item Problem formulation: Draft optimisation problem specified with decision variables, objective function (ELBO), and constraints

\item Mathematical derivations\footnote{Desperiately in need of some kind of peer review}:   
\begin{itemize}
\item Draft ELBO derivation from first principles
\item Draft CAVI algorithm with closed-form updates
\item Analytical gradient expressions for all 11 variational parameters (unverified) 
\item Potential Hessian block computations for Newton's method
\item Attempt BFGS quasi-Newton formulation
\item Attempt at Penalty method formulation for constrained optimisation
\end{itemize}

\item Validation strategy: Established CAVI as analytical benchmark for correctness checking
\end{enumerate}

\subsection*{Preliminary Results}

Implementation is pending, the mathematical analysis hints\footnote{I havent been over this enough to make any claim of completeness}
 
\newpage
\subsection*{Next Steps}

\begin{enumerate}
\item Implement CAVI baseline in Python (Week 7)
\item Implement parameter ascent with line search (Week 7)
\item Implement Newton's method with regularisation (Week 8)
\item Implement BFGS quasi-Newton (Week 8)
\item Monte Carlo comparison: convergence iterations, wall-clock time, ELBO values (Week 8-9)
\item Final report with performance analysis (Week 12)
\end{enumerate}

\section*{Alignment with Course Requirements}

  

 The VI problem is meant to represent  real computational challenges, making the comparison of optimisation methods   valuable,  relevant, and interesting.

\section*{Timeline}

TBD

\section*{Summary}

I've set up the optimisation problem and derived the math. The next step is to code four algorithms---CAVI, gradient ascent, Newton, BFGS---and see which one wins, and why.

\end{document}


